% ==============================================================================
% SECTION II: SYSTEM MODEL, THREAT MODEL & THE ANOMALY DETECTION TRILEMMA
% Target Venues: IEEE S&P / ACM CCS / USENIX Security / NDSS
% ==============================================================================

\section{System Model, Threat Model, and Defensible Research Gap}
\label{sec:threat_model}

In this section, we establish the mathematical system model for decentralized multi-tenant IoT edge computing, formalize the dual-spectrum adversary capabilities, and articulate the fundamental research gap through the lens of the Anomaly Detection Trilemma.

\subsection{Decentralized Multi-Tenant IoT Edge System Model}
\label{subsec:system_model}

We model an enterprise-scale, decentralized IoT infrastructure consisting of $M$ edge gateways $\mathcal{G} = \{\mathcal{G}_1, \mathcal{G}_2, \dots, \mathcal{G}_M\}$ interconnected over an encrypted wide-area network (WAN) with an orchestration parameter server $\mathcal{S}$. Each gateway $\mathcal{G}_i$ resides at the boundary of a dedicated physical or administrative subnet, monitoring real-time network flow telemetry generated by specialized IoT end-devices (e.g., smart power meters, industrial PLCs, surveillance cameras, or environmental sensors). 

Let $\mathcal{X} \subseteq \mathbb{R}^D$ denote the continuous $D$-dimensional feature space of network flow statistics (e.g., inter-arrival times, payload byte distributions, TCP window metrics, packet length moments). Due to application-specific operational protocols, nominal traffic observed by gateway $\mathcal{G}_i$ is supported on a compact, low-dimensional Riemannian sub-manifold $\mathcal{M}_i \subset \mathbb{R}^D$ with intrinsic dimension $r_i \ll D$:
\begin{equation}
\label{eq:nominal_manifold}
\mathcal{M}_i \triangleq \left\{ x \in \mathbb{R}^D \;\middle|\; \norm{(I - U_i U_i^\top)(x - \mu_i)}_2 \le \tau_i, \; \mathcal{D}_{\text{Maha}}^2(x; \mu_i, \Lambda_i) \le \chi^2_{r_i}(1-\alpha) \right\},
\end{equation}
where $\mu_i = \mathbb{E}_{x \sim \mathcal{P}_i}[x] \in \mathbb{R}^D$ is the empirical centroid of gateway $\mathcal{G}_i$, $U_i \in \mathbb{R}^{D \times r_i}$ is the orthonormal projection basis spanning the top-$r_i$ principal eigenspace of nominal covariance $\Sigma_i$, $\Lambda_i \in \mathbb{R}^{r_i \times r_i}$ is the diagonal eigenvalue matrix, $\tau_i$ denotes the orthogonal reconstruction tolerance, and $\mathcal{D}_{\text{Maha}}^2(x; \mu_i, \Lambda_i) \triangleq (x - \mu_i)^\top U_i \Lambda_i^{-1} U_i^\top (x - \mu_i)$ denotes the intra-subspace Mahalanobis distance bounded by the chi-squared distribution at significance level $1-\alpha$.

\subsubsection{Extreme Non-IID Manifold Disjointness}
In a multi-tenant edge deployment, different edge gateways observe fundamentally disparate application domains (e.g., high-frequency SCADA polling at $\mathcal{G}_1$ vs. bursty video streaming at $\mathcal{G}_2$). Consequently, the marginal data distributions $\mathcal{P}_i \neq \mathcal{P}_j$ exhibit severe Non-IID skew, and their nominal support manifolds are mutually disjoint in the ambient space:
\begin{equation}
\label{eq:disjoint_manifolds}
\mathcal{M}_i \cap \mathcal{M}_j = \emptyset \quad \forall i \neq j, \quad \text{with} \quad \inf_{x_i \in \mathcal{M}_i, x_j \in \mathcal{M}_j} \norm{x_i - x_j}_2 \ge \Delta_{\min} > 0,
\end{equation}
where $\Delta_{\min}$ characterizes the minimum manifold separation distance. Under stringent data confidentiality policies (e.g., GDPR, HIPAA), gateways are strictly prohibited from transmitting raw network telemetry $x$ to the parameter server $\mathcal{S}$ or to peer gateways.

\subsection{Threat Model and Adversary Capabilities}
\label{subsec:threat_model}

We adopt the standard threat modeling conventions of top-tier security venues (IEEE S\&P, ACM CCS), characterizing both the adversary's operational objectives and environmental trust assumptions:

\subsubsection{Adversary Objectives \& Dual-Spectrum Attack Taxonomy}
We consider an external or compromised internal network adversary $\mathcal{A}$ possessing the capability to inject malicious traffic flows into the ingress/egress links of any monitored subnet. The adversary's objective is to execute malicious actions without triggering detection alarms at the local gateways. We classify malicious traffic into a dual-spectrum taxonomy:
\begin{enumerate}
    \item \textbf{Macroscopic Volumetric Attacks}: These encompass massive distributed denial-of-service (DDoS) floods (e.g., UDP/SYN floods, HTTP GET floods, Botnet amplification) generated by automated botnets such as Mirai and Bashlite. Volumetric attacks induce dramatic surges in packet count, byte volume, and flow frequency, projecting anomalous feature vectors $x_{\text{att}}$ deep into out-of-distribution space ($\norm{x_{\text{att}} - \mu_i}_2 \gg \text{diam}(\mathcal{M}_i)$).
    \item \textbf{Microscopic Stealthy \& Structural Attacks}: These encompass multi-stage targeted intrusions, such as slow horizontal port scans, asynchronous command-and-control (C2) beaconing, credential brute-forcing, and lateral pivoting. Sophisticated adversaries intentionally throttle transmission rates and randomize flow timing to mimic normal marginal feature distributions. Consequently, individual feature values for $x_{\text{att}}$ fall within the allowable range of nominal traffic, and the anomaly manifests exclusively through abnormal relational topology and spatial distance transitions among neighboring points ($D(x_{\text{att}}, \mathcal{N}_k)$).
\end{enumerate}

\subsubsection{Trust Assumptions \& Threat Boundary}
We assume that the edge gateway hardware, operating system, and sensor capture interfaces (e.g., DPDK/eBPF drivers) are trusted and execute within a secure hardware root-of-trust (e.g., ARM TrustZone). The parameter server $\mathcal{S}$ is modeled as \emph{honest-but-curious}: it faithfully executes the federated aggregation protocol, but may attempt to infer private gateway traffic characteristics from transmitted gradient vectors or model updates. We consider an uncoordinated adversarial environment where edge gateways may be subjected to independent, asynchronous zero-day attacks without prior attack signatures or malicious labels available during training.

\subsection{The Anomaly Detection Trilemma: Gap vs. Baselines}
\label{subsec:trilemma}

To rigorously articulate why existing intrusion detection paradigms fail in decentralized IoT environments, we formulate the \textbf{Anomaly Detection Trilemma for IoT Edge Networks}. An ideal edge NIDS must simultaneously satisfy three fundamental properties:
\begin{itemize}
    \item \textbf{P1 (Stealthy Relational Sensitivity)}: High detection efficacy against microscopic, structural attacks that mimic marginal distributions.
    \item \textbf{P2 (Non-IID Drift Robustness)}: Resilience against extreme client manifold skew ($\mathcal{M}_i \cap \mathcal{M}_j = \emptyset$) without representation collapse or gradient cancellation.
    \item \textbf{P3 (Edge Resource Viability)}: Sub-millisecond streaming line-rate inference with bounded, lightweight memory footprints ($<100$\,MB RAM).
\end{itemize}

As summarized in Table~\ref{tab:trilemma_comparison}, existing paradigms sacrifice at least one fundamental property:

\begin{table*}[t]
\centering
\caption{The Anomaly Detection Trilemma: Comparative analysis of detection paradigms in multi-tenant Non-IID edge networks.}
\label{tab:trilemma_comparison}
\adjustbox{max width=\textwidth}{
\begin{tabular}{lcccccc}
\toprule
\textbf{Detection Paradigm} & \textbf{Representative Models} & \textbf{Detection Mechanism} & \textbf{P1: Stealthy Sensitivity} & \textbf{P2: Non-IID Robustness} & \textbf{P3: Edge Resource Viability} & \textbf{Catastrophic Failure Mode} \\
\midrule
\textbf{Signature / Rule-based} & Snort~\cite{roesch1999snort}, Zeek~\cite{paxson1999bro} & Deterministic Pattern Match & \xmark~Low & \cmark~High (Static Rules) & \cmark~Line-rate ($<0.01$\,ms) & Complete Blindness to Zero-Day Attacks \\
\textbf{Reconstruction / Statistical} & Kitsune~\cite{mirsky2018kitsune}, Fed-AE~\cite{sakurada2014anomaly} & Residual $\norm{x - D(E(x))}_2^2$ & \xmark~Low (Shortcut Trap) & \cmark~Moderate & \cmark~Fast ($0.002$\,ms, $18.7$\,MB) & Reconstruction Shortcut on Stealthy C2 \\
\textbf{Tree / Density Ensembles} & Isolation Forest~\cite{liu2008isolation}, OCSVM~\cite{scholkopf2001estimating} & Axis-Aligned Space Slicing & \xmark~Moderate & \xmark~Poor (FL Incompatible) & \cmark~Low Memory & Axis-Aligned Bias, Massive False Alarms \\
\textbf{Spectral / Analytical Bounds} & LOC-NFST~\cite{nguyen2024locnfst}, KNFST~\cite{bodesheim2013kernel} & Exact Null Space ($S_w w = 0$) & \cmark~High & \xmark~Poor (Drift Erosion) & \xmark~Prohibitive ($>500$\,MB RAM) & Null-Space Erosion \& $\mathcal{O}(N^2)$ RAM Scale \\
\textbf{Naive Distance-Ranking Graph} & Naive Fed-LUNAR~\cite{goodge2022lunar} & Topological $k$-NN Ranking & \cmark~High & \xmark~Fatal (Cancel/Inversion) & \cmark~Light ($50$\,MB) & Distance Inversion ($0.15\%$) \& Gradient Clash \\
\midrule
\textbf{Fed-LUNAR (Proposed)} & \textbf{MSSP + FSDS + CMNP + DROGA} & \textbf{Orthogonal Aligned $k$-NN} & \textbf{\cmark~Superior ($99.4\%$ F1)} & \textbf{\cmark~Immune ($\alpha=0.1$)} & \textbf{\cmark~Line-rate ($0.001$\,ms, $54$\,MB)} & \textbf{None (Overcomes Inversion \& Clash)} \\
\bottomrule
\end{tabular}
}
\end{table*}

\subsubsection{The Reconstruction Shortcut Trap in Autoencoders}
Autoencoder-based NIDS (e.g., Kitsune~\cite{mirsky2018kitsune}, FedAutoEncoder~\cite{sakurada2014anomaly}) minimize empirical reconstruction loss $\mathcal{L}_{\text{AE}} = \frac{1}{N}\sum_{i=1}^N \norm{x_i - \mathcal{D}(\mathcal{E}(x_i))}_2^2$. Because deep neural networks are universal function approximators with smooth inductive bias, the latent bottleneck $\mathcal{E}(x)$ learns continuous manifold charts that span unpopulated low-density regions. When the adversary injects slow, structural attack packets whose marginal feature distributions intersect benign telemetry, the decoder reconstructs the anomalous packets with negligible residual error ($\norm{x_{\text{att}} - \hat{x}_{\text{att}}}_2^2 \le \tau_{\text{threshold}}$). This phenomenon, known as \emph{reconstruction shortcutting}~\cite{gong2019memorizing}, produces unacceptable false-negative rates on stealthy botnets.

\subsubsection{Catastrophic Null-Space Erosion in Analytical Spectral Methods}
Analytical methods, such as the Kernel Null Foley-Sammon Transform (KNFST)~\cite{bodesheim2013kernel} and LOC-NFST~\cite{nguyen2024locnfst}, map training samples into an exact null space satisfying $w^\top S_w w = 0$ while maximizing between-class scatter $w^\top S_b w > 0$. While closed-form solutions achieve instantaneous one-round aggregation, they suffer two fatal bottlenecks in continuous streaming environments:
\begin{enumerate}
    \item \textbf{Erosion of Null Space under Non-IID Streaming Drift}: In streaming network topologies where flow volumes rapidly scale ($N \gg D$), empirical covariance matrices become strictly positive definite ($S_w \succ 0$), eliminating the zero-eigenvalue null space. Non-stationary streaming drift forces iterative rank-one updates that rapidly erode the projection subspace, causing detection thresholds to degenerate ($\tau \to 0$).
    \item \textbf{Prohibitive Memory Footprint}: Computing and updating kernel Gram matrices $\mathbf{K} \in \mathbb{R}^{N \times N}$ requires $\mathcal{O}(N^2)$ memory storage. As empirically demonstrated in Section~\ref{sec:experiments}, LOC-NFST consumes $434.8$\,MB to $958.2$\,MB of active RAM, rendering it completely unviable for low-cost edge gateways (e.g., IoT routers with 512\,MB total memory).
\end{enumerate}

\subsubsection{The Unique Advantage of Graph Distance-Ranking & The Unsolved Dilemma}
Graph-based distance-ranking architectures (LUNAR~\cite{goodge2022lunar}) circumvent the reconstruction shortcut by abandoning pixel/feature-wise reconstruction in favor of learning local density gradients from sorted neighbor distance vectors $d(x) = [d_1(x), \dots, d_k(x)]^\top$. However, as introduced in Section~\ref{sec:intro}, attempting to deploy distance-ranking in a decentralized federated network triggers two catastrophic vulnerabilities: OOD Distance Inversion (collapsing AUC to $0.15\%$) and Cross-Manifold Negative Gradient Cancellation (destroying federated convergence). Resolving these two theoretical dilemmas while preserving sub-millisecond edge streaming throughput constitutes the core objective of Fed-LUNAR.
